%% Notes
%% This document requires XeLaTeX or LuaLaTeX to compile due to the use of system fonts.
%% Remove the section named "Remove for vanilla LaTeX" if you want to compile with pdfLaTeX.
%%

\documentclass[a4paper,12pt]{article}

%% Custom fonts - Remove for vanilla LaTeX
\usepackage{fontspec}
\setmainfont{RobotoSerif-Regular.ttf}[
    Path = Assets/Fonts/Roboto_Serif/,
    Extension = .ttf,
    BoldFont = RobotoSerif-Bold.ttf,
    ItalicFont = RobotoSerif-Italic.ttf,
    BoldItalicFont = RobotoSerif-BoldItalic.ttf
]
%% /Custom fonts - Remove for vanilla LaTeX
%% Add this if not using XeLaTeX or LuaLaTeX
%\usepackage[utf8]{inputenc}

\usepackage[none]{hyphenat} %% Ensures no hyphenation
\emergencystretch=2em % Alows more empty space between words on lines
\usepackage[final]{microtype} % Improves text appearance - apparently ???

\usepackage{graphics} % for \scalebox
\usepackage{amsmath} % Advanced math typesetting
\usepackage{amssymb} % for math symbols
\usepackage{float} % Improved interface for floating objects
\usepackage[left=1in,right=1in,top=1in,bottom=1in]{geometry} % Smaller margins


\usepackage{tikz} % Advanced drawings
\usetikzlibrary{arrows.meta, angles, quotes}


%% Remove numbering and add symbol before section and subsection titles
\usepackage{titlesec}
\titleformat{\section}
    {\normalfont\Large\bfseries}
    {\thesection \quad}
    {0pt}
    {}
\titleformat{\subsection}
    {\normalfont\large\bfseries}
    {\thesubsection \hspace{0.05em} $-$}
    {0.25em}
    {}
\titleformat{\subsubsection}
    {\normalfont\normalsize\bfseries}
    {}
    {0em}
    {}
\titleformat{\paragraph}
    [runin]
    {\normalfont\normalsize\bfseries\itshape}
    {}
    {0em}
    {}
%% /Remove numbering and add symbol before section and subsection titles

%% Reduce page-break inside sections
\usepackage{needspace}

% discourage page breaks inside paragraphs
\interlinepenalty=10000
\clubpenalty=10000
\widowpenalty=10000
%% /Reduce page-break inside sections.

\setlength{\parindent}{0pt} % No paragraph indentation

%% header and footer
\usepackage{fancyhdr}
\fancyhf{}
\fancyhead[L]{\leftmark}
\fancyhead[R]{\thepage}
\fancyfoot{}
\renewcommand{\sectionmark}[1]{\markboth{#1}{}}
\setlength{\headheight}{14.5pt}
%% /header and footer

% For multiplication table
\usepackage{array}
\usepackage{multirow}
\usepackage{colortbl}
\definecolor{Gray}{gray}{0.85}
\newcolumntype{a}{>{\columncolor{Gray}}c}

\usepackage{tabularray}
\usepackage[x11names]{xcolor}
% /For multiplication table

% For bemærk boxes
\usepackage{xcolor}
\usepackage{tcolorbox}
\tcbuselibrary{breakable}
% /For bemærk boxes

% For pretty tables
\usepackage{booktabs}
% /For pretty tables

% For compact TOC

\usepackage{titletoc}
\titlecontents{section}               % Section
    [0em]                             % Left
    {\vspace*{\baselineskip}}         % Above code
    {\bfseries\thecontentslabel\quad} % Numbered format
    {\bfseries\thecontentslabel}      % Numberless format
    { \textit{S. \thecontentspage}}         % Filler
    []                                % Separator
\titlecontents*{subsection}           % Section
    [0em]                             % Left
    {\space}                          % Above code
    {\thecontentslabel~}              % Numbered format
    {}                                % Numberless format
    {}                                % Filler
    [\hspace{2em}]                    % Separator

\usepackage[hidelinks]{hyperref}
% /For compact TOC